% GENERATED FILE. Edit canonical lessons, then run npm run generate:books.
% canonical-source-hash: fnv1a64:412dcf798699e208
% canonical-lessons: PA-C291-thandha, PA-C291-niggha, PA-C291-tatta, PA-C291-kosa, PA-C291-thaarna

\chapter{Warm and Cool}
\label{ch:pa-a2-warm-and-cool}
\hlchaptermodality{291}

\begin{chapteropening}
\textbf{By the end of this chapter:} \emph{I can say how warm or cool something feels.}

Complete the last lesson of chapter 291: I can say how warm or cool something feels.
\end{chapteropening}

\section[\texorpdfstring{\pa{ਠੰਢਾ} (ṭhaṇḍhā)}{ਠੰਢਾ (ṭhaṇḍhā)}]{\pa{ਠੰਢਾ} (ṭhaṇḍhā) — cold; cool}

\label{lesson:PA-C291-thandha}



\begin{quote}
Before the new one: say the Punjabi for to go down, then the Punjabi for to blow.
\end{quote}

\subsection*{You'll want to know: \pa{ਠੰਢਾ}}
\textbf{\pa{ਠੰਢਾ}} — \emph{ṭhaṇḍhā} — \textquotedblleft{}cold; cool\textquotedblright{}.

\begin{quote}
\pa{ਘੜੇ} \pa{ਦਾ} \pa{ਪਾਣੀ} \pa{ਠੰਢਾ} \pa{ਹੁੰਦਾ} \pa{ਹੈ।}
\end{quote}

\begin{quote}
\emph{ghaṛe dā pāṇī ṭhaṇḍhā hundā hai.}
\end{quote}

\begin{quote}
Water from a clay pot is cool.
\end{quote}

\begin{cousinweb}
\textbf{\pa{ਠੰਢਾ}} (\emph{ṭhaṇḍhā}) goes back to Sanskrit \emph{stabdha}, \textquotedblleft{}stiff, solidified (as water)\textquotedblright{}; CDIAL lists Punjabi \emph{ṭhāḍhā} under it only as \textquotedblleft{}upright\textquotedblright{}, a word it marks as taken from Hindi; the sense \textquotedblleft{}cold\textquotedblright{} is the one it gives for Lahnda \emph{ṭhaḍḍhā} and Hindi \emph{ṭhaṇḍhā} (CDIAL 13676).
\end{cousinweb}

\begin{grammarlens}[title={What you've built}]
One more word for this chapter.
\end{grammarlens}

\subsection*{Your turn}
\begin{itemize}
\raggedright
  \item \emph{Say it:} \emph{ṭhaṇḍhā}
  \item \emph{Say it:} \emph{ṭhaṇḍhā}, once more
  \item \emph{Say it:} \emph{ghaṛe dā pāṇī ṭhaṇḍhā hundā hai.} — Water from a clay pot is cool.
  \item \emph{Recall:} say the Punjabi for to go down, then the Punjabi for to blow, then say \emph{ṭhaṇḍhā} again
\end{itemize}

\begin{tcolorbox}[breakable,colback=teal!4,colframe=teal!35!black,title={Before you move on}]
What does \pa{ਠੰਢਾ} mean? (\textquotedblleft{}Cold; cool\textquotedblright{}.) Say it once more.
\end{tcolorbox}
\section[\texorpdfstring{\pa{ਨਿੱਘਾ} (nigghā)}{ਨਿੱਘਾ (nigghā)}]{\pa{ਨਿੱਘਾ} (nigghā) — pleasantly warm; cosy}

\label{lesson:PA-C291-niggha}



\begin{quote}
Before the new one: say the Punjabi for to blow, then the Punjabi for cool.
\end{quote}

\subsection*{You'll want to know: \pa{ਨਿੱਘਾ}}
\textbf{\pa{ਨਿੱਘਾ}} — \emph{nigghā} — \textquotedblleft{}pleasantly warm; cosy\textquotedblright{}.

\begin{quote}
\pa{ਇਹ} \pa{ਕੰਬਲ} \pa{ਬਹੁਤ} \pa{ਨਿੱਘਾ} \pa{ਹੈ।}
\end{quote}

\begin{quote}
\emph{eh kambal bahut nigghā hai.}
\end{quote}

\begin{quote}
This blanket is very warm.
\end{quote}

\begin{cousinweb}
\textbf{\pa{ਨਿੱਘਾ}} (\emph{nigghā}) is inherited from a reconstructed \emph{*nidāghya}, \textquotedblleft{}heat\textquotedblright{}; CDIAL lists Punjabi \emph{nigghā}, \textquotedblleft{}moderately warm\textquotedblright{}, under it (CDIAL 7195).
\end{cousinweb}

\begin{grammarlens}[title={What you've built}]
One more word for this chapter.
\end{grammarlens}

\subsection*{Your turn}
\begin{itemize}
\raggedright
  \item \emph{Say it:} \emph{nigghā}
  \item \emph{Say it:} \emph{nigghā}, once more
  \item \emph{Say it:} \emph{eh kambal bahut nigghā hai.} — This blanket is very warm.
  \item \emph{Recall:} say the Punjabi for to blow, then the Punjabi for cool, then say \emph{nigghā} again
  \item \emph{Write it:} copy \textbf{\pa{ਨਿੱਘਾ}} once, checking the whole word
\end{itemize}

\begin{tcolorbox}[breakable,colback=teal!4,colframe=teal!35!black,title={Before you move on}]
What does \pa{ਨਿੱਘਾ} mean? (\textquotedblleft{}Pleasantly warm; cosy\textquotedblright{}.) Say it once more.
\end{tcolorbox}
\section[\texorpdfstring{\pa{ਤੱਤਾ} (tattā)}{ਤੱਤਾ (tattā)}]{\pa{ਤੱਤਾ} (tattā) — hot to the touch}

\label{lesson:PA-C291-tatta}



\begin{quote}
Before the new one: say the Punjabi for cool, then the Punjabi for pleasantly warm.
\end{quote}

\subsection*{You'll want to know: \pa{ਤੱਤਾ}}
\textbf{\pa{ਤੱਤਾ}} — \emph{tattā} — \textquotedblleft{}hot to the touch\textquotedblright{}.

\begin{quote}
\pa{ਚਾਹ} \pa{ਬਹੁਤ} \pa{ਤੱਤੀ} \pa{ਹੈ।}
\end{quote}

\begin{quote}
\emph{chāh bahut tattī hai.}
\end{quote}

\begin{quote}
The tea is very hot.
\end{quote}

\begin{cousinweb}
\textbf{\pa{ਤੱਤਾ}} (\emph{tattā}) is inherited from Sanskrit \emph{tapta}, \textquotedblleft{}heated, hot\textquotedblright{}; CDIAL lists Punjabi \emph{tattā}, \textquotedblleft{}hot, burning\textquotedblright{}, under it (CDIAL 5679).
\end{cousinweb}

\begin{grammarlens}[title={What you've built}]
One more word for this chapter.
\end{grammarlens}

\subsection*{Your turn}
\begin{itemize}
\raggedright
  \item \emph{Say it:} \emph{tattā}
  \item \emph{Say it:} \emph{tattā}, once more
  \item \emph{Say it:} \emph{chāh bahut tattī hai.} — The tea is very hot.
  \item \emph{Recall:} say the Punjabi for cool, then the Punjabi for pleasantly warm, then say \emph{tattā} again
  \item \emph{Write it:} copy \textbf{\pa{ਤੱਤਾ}} once, checking the whole word
\end{itemize}

\begin{tcolorbox}[breakable,colback=teal!4,colframe=teal!35!black,title={Before you move on}]
What does \pa{ਤੱਤਾ} mean? (\textquotedblleft{}Hot to the touch\textquotedblright{}.) Say it once more.
\end{tcolorbox}
\section[\texorpdfstring{\pa{ਕੋਸਾ} (kosā)}{ਕੋਸਾ (kosā)}]{\pa{ਕੋਸਾ} (kosā) — lukewarm}

\label{lesson:PA-C291-kosa}



\begin{quote}
Before the new one: say the Punjabi for pleasantly warm, then the Punjabi for hot to the touch.
\end{quote}

\subsection*{You'll want to know: \pa{ਕੋਸਾ}}
\textbf{\pa{ਕੋਸਾ}} — \emph{kosā} — \textquotedblleft{}lukewarm\textquotedblright{}.

\begin{quote}
\pa{ਮੈਂ} \pa{ਕੋਸੇ} \pa{ਪਾਣੀ} \pa{ਨਾਲ} \pa{ਨਹਾਉਂਦਾ} \pa{ਹਾਂ।}
\end{quote}

\begin{quote}
\emph{maiṁ kose pāṇī nāl nahāundā hāṁ.}
\end{quote}

\begin{quote}
I bathe with lukewarm water.
\end{quote}

\begin{cousinweb}
\textbf{\pa{ਕੋਸਾ}} (\emph{kosā}) is inherited from a reconstructed \emph{*kōṣma}, \textquotedblleft{}warm\textquotedblright{}; CDIAL lists Punjabi \emph{kōssā}, \textquotedblleft{}lukewarm\textquotedblright{}, under it (CDIAL 3552).
\end{cousinweb}

\begin{grammarlens}[title={What you've built}]
One more word for this chapter.
\end{grammarlens}

\subsection*{Your turn}
\begin{itemize}
\raggedright
  \item \emph{Say it:} \emph{kosā}
  \item \emph{Say it:} \emph{kosā}, once more
  \item \emph{Say it:} \emph{maiṁ kose pāṇī nāl nahāundā hāṁ.} — I bathe with lukewarm water.
  \item \emph{Recall:} say the Punjabi for pleasantly warm, then the Punjabi for hot to the touch, then say \emph{kosā} again
  \item \emph{Write it:} copy \textbf{\pa{ਕੋਸਾ}} once, checking the whole word
\end{itemize}

\begin{tcolorbox}[breakable,colback=teal!4,colframe=teal!35!black,title={Before you move on}]
What does \pa{ਕੋਸਾ} mean? (\textquotedblleft{}Lukewarm\textquotedblright{}.) Say it once more.
\end{tcolorbox}
\section[\texorpdfstring{\pa{ਠਾਰਨਾ} (ṭhārnā)}{ਠਾਰਨਾ (ṭhārnā)}]{\pa{ਠਾਰਨਾ} (ṭhārnā) — to cool something down}

\label{lesson:PA-C291-thaarna}



\begin{quote}
Before the new one: say the Punjabi for hot to the touch, then the Punjabi for lukewarm.
\end{quote}

\subsection*{You'll want to know: \pa{ਠਾਰਨਾ}}
\textbf{\pa{ਠਾਰਨਾ}} — \emph{ṭhārnā} — \textquotedblleft{}to cool something down\textquotedblright{}.

\begin{quote}
\pa{ਮੈਂ} \pa{ਬੱਚੇ} \pa{ਲਈ} \pa{ਦੁੱਧ} \pa{ਠਾਰਦਾ} \pa{ਹਾਂ।}
\end{quote}

\begin{quote}
\emph{maiṁ bachche laī dudh ṭhārdā hāṁ.}
\end{quote}

\begin{quote}
I cool the milk down for the child.
\end{quote}

Past (I …): \pa{ਮੈਂ} \pa{ਦੁੱਧ} \pa{ਠਾਰਿਆ}, \emph{maiṁ dudh ṭhāriā}.

\begin{cousinweb}
\textbf{\pa{ਠਾਰਨਾ}} (\emph{ṭhārnā}) is the causative partner of \emph{ṭharnā}, inherited from a reconstructed \emph{*sthārayati}, \textquotedblleft{}makes firm\textquotedblright{}; CDIAL lists Punjabi \emph{ṭhārnā}, \textquotedblleft{}to cool\textquotedblright{}, under it (CDIAL 13742).
\end{cousinweb}

\begin{grammarlens}[title={What you've built}]
One more word for this chapter.
\end{grammarlens}

\subsection*{Your turn}
\begin{itemize}
\raggedright
  \item \emph{Say it:} \emph{ṭhārnā}
  \item \emph{Say it:} \emph{ṭhārnā}, once more
  \item \emph{Say it:} \emph{maiṁ bachche laī dudh ṭhārdā hāṁ.} — I cool the milk down for the child.
  \item \emph{Recall:} say the Punjabi for hot to the touch, then the Punjabi for lukewarm, then say \emph{ṭhārnā} again
  \item \emph{Write it:} copy \textbf{\pa{ਠਾਰਨਾ}} once, checking the whole word
\end{itemize}

\begin{tcolorbox}[breakable,colback=teal!4,colframe=teal!35!black,title={Before you move on}]
What does \pa{ਠਾਰਨਾ} mean? (\textquotedblleft{}To cool something down\textquotedblright{}.) Say it once more.
\end{tcolorbox}
